\documentclass[a4paper,11pt]{article}

\usepackage[a4paper,margin=1.8cm]{geometry}
\usepackage{fontspec}
\usepackage{polyglossia}
\setmainlanguage{russian}
\setmainfont{Arial}
\setmonofont{Arial}

\usepackage{xcolor}
\usepackage{array}
\usepackage{enumitem}
\usepackage{titlesec}
\usepackage{hyperref}
\usepackage{fancyhdr}

\definecolor{accent}{HTML}{2457A6}
\definecolor{darkgray}{HTML}{39424E}
\definecolor{noteBackground}{HTML}{FFF4CC}
\definecolor{noteBorder}{HTML}{D18B00}

\hypersetup{
  colorlinks=true,
  linkcolor=accent,
  urlcolor=accent,
  pdftitle={Sprint 1 PMLDL}
}

\titleformat{\section}
  {\Large\bfseries\color{accent}}
  {}{0pt}{}

\pagestyle{fancy}
\fancyhf{}
\setlength{\headheight}{14pt}
\lhead{PMLDL · LLM Classification Fine-tuning}
\rhead{Sprint 1}
\cfoot{\thepage}
\renewcommand{\headrulewidth}{0.4pt}

\newcommand{\code}[1]{\texttt{#1}}
\newcommand{\experiment}[2]{\paragraph{\textcolor{accent}{#1 --- #2}}}

\begin{document}

\begin{center}
  {\Huge\bfseries\color{accent} Sprint 1 PMLDL}\\[0.4em]
  {\large Распределение baseline и экспериментов команды}
\end{center}

\section*{Распределение заданий}

\renewcommand{\arraystretch}{1.35}
\begin{center}
\begin{tabular}{|>{\raggedright\arraybackslash}p{0.23\textwidth}|
                >{\raggedright\arraybackslash}p{0.34\textwidth}|
                >{\raggedright\arraybackslash}p{0.34\textwidth}|}
\hline
\textbf{Участник} & \textbf{Основные задания} & \textbf{Дополнительные задания} \\
\hline
Данил & E9 --- ensemble/blending & E8, E10 \\
\hline
Данил Куджи & E1 --- pretrained transformer & E4 \\
\hline
Карим & E2 --- LoRA/QLoRA и fine-tuning & E6 \\
\hline
Аслан & E3 --- A/B-симметричная архитектура & E5 \\
\hline
Никита & B0 --- class-prior baseline; B1 --- structural baseline & --- \\
\hline
Арсений & B2 --- TF-IDF baseline & B3, E7 \\
\hline
\end{tabular}
\end{center}

\vspace{0.5em}
\noindent\fcolorbox{noteBorder}{noteBackground}{\parbox{0.94\textwidth}{
  \textbf{Важно:} дополнительные задания необязательны. Их можно выполнить
  после следующего stage, если останутся время и ресурсы. В Sprint 1 приоритет
  имеют основные задания и baseline B0--B3.
}}

\section*{Описание baseline и экспериментов}

\experiment{B0}{Class-prior baseline}
Предсказывать самый частый класс или использовать распределение классов,
рассчитанное на обучающей выборке. Это минимальная контрольная точка: любой
рабочий подход должен превосходить её по \code{log\_loss}.

\experiment{B1}{Structural baseline}
Использовать длины prompt и ответов A/B, количество слов и токенов,
соотношение длин, специальные символы и признаки пустых ответов. Поверх этих
признаков обучить Logistic Regression и проверить, насколько простая структура
диалога уже объясняет предпочтение ответа.

\experiment{B2}{TF-IDF baseline}
Преобразовать prompt и ответы A/B в word-level и character-level TF-IDF-признаки.
Обучить Logistic Regression или Linear SVM с вероятностным выводом. Этот
baseline нужен как сильное и быстрое классическое решение для сравнения
с нейронными моделями.

\experiment{B3}{Frozen embeddings}
Получить готовые sentence embeddings с помощью pretrained encoder без его
дообучения. Поверх эмбеддингов обучить Logistic Regression или небольшой MLP.
Так проверяется, достаточно ли универсальных представлений текста без адаптации
энкодера под задачу соревнования.

\experiment{E1}{Full transformer fine-tuning}
Дообучить pretrained-модель на классификации трёх классов: победа модели A,
победа модели B или tie. Можно выбрать DeBERTa, RoBERTa, ModernBERT или
аналогичную модель. Сравнивать нужно качество и ресурс обучения относительно
лучшего baseline.

\experiment{E2}{LoRA/QLoRA fine-tuning}
Сравнить полное дообучение с LoRA и QLoRA на одной и той же pretrained-модели.
Исследовать rank, dropout, learning rate и количество обучаемых параметров,
сохраняя остальные условия фиксированными.

\experiment{E3}{A/B-симметричная архитектура}
Пропускать ответы A и B через общий encoder отдельно. Использовать признаки
\code{h\_A-h\_B}, \code{|h\_A-h\_B|} и признаки prompt. Проверить, повышает ли
такая конструкция устойчивость модели к перестановке ответов.

\experiment{E4}{Swap augmentation и consistency loss}
Создавать копию каждого примера с переставленными ответами A и B и менять target.
Дополнительно штрафовать модель за несогласованные предсказания после перестановки.
Сравнить обучение без augmentation, только с augmentation и с consistency loss.

\experiment{E5}{Стратегия truncation}
Сравнить head-only, tail-only, head-tail и turn-aware truncation при одинаковой
модели и бюджете токенов. Проверить, какая стратегия лучше сохраняет важную
информацию в длинных диалогах.

\experiment{E6}{Функция потерь}
Сравнить обычную cross-entropy, label smoothing, class weights и focal loss.
Цель --- проверить, помогает ли изменение loss при дисбалансе классов и как оно
влияет на calibration и итоговый \code{log\_loss}.

\experiment{E7}{Калибровка вероятностей}
Применить temperature scaling или vector scaling после обучения, не меняя
порядок классов и сами предсказания модели. Оценивать изменение качества
вероятностей по основной метрике \code{log\_loss}.

\experiment{E8}{Test-time augmentation}
Использовать несколько вариантов truncation или окон одного диалога во время
inference и усреднять итоговые вероятности. Сравнить прирост качества с
дополнительным временем инференса.

\experiment{E9}{Blending моделей}
Смешать вероятности TF-IDF, embeddings и transformer-моделей, которые допускают
разные типы ошибок. Сравнить простое усреднение с подбором весов на validation.

\experiment{E10}{Устойчивость к seed}
Запустить одну конфигурацию с несколькими seed, усреднить вероятности и оценить
средний результат и разброс качества. Это показывает, насколько результат
воспроизводим и не зависит от случайной инициализации.

\vfill
\begin{center}
  \small\color{darkgray}
  Sprint 1 считается выполненным после получения основных baseline и экспериментов.
\end{center}

\end{document}
